% 张三慧《大学物理学》第4版第二卷第28章原书“思考题”28.1--28.14
% 主教材扫描PDF物理页423（印刷406），单独类别，不占章末习题28.1--28.34号。
\section*{第28章\quad 原子中的电子：思考题}
\paragraph{思考题 28.1}
为什么用确定的经典轨道描述原子内部电子的运动状态是不恰当的？
\paragraph{思考题 28.2}
能级“简并”是什么意思？忽略电子自旋时，氢原子第 \(n\) 能级在中心库仑场中由多少种不同的量子态构成？
\paragraph{思考题 28.3}
钾原子的价电子能级主要由哪些量子数决定？为什么不能将单电子氢原子的“仅依赖 \(n\)”简单推广到多电子原子？
\paragraph{思考题 28.4}
反氢原子由反质子及围绕它的正电子组成，它在相同物理环境下的光谱与氢光谱是否应当相同？说明所依赖的对称性前提。
\paragraph{思考题 28.5}
主量子数 \(n=3\) 的壳层有多少次壳层？每次壳层最多容纳多少个电子？
\paragraph{思考题 28.6}
用经典电子圆轨道模型推出氢原子能量
\(E_{\rm class}=-e^2/(8\pi\epsilon_0r)\)，
并与Bohr能级比较，说明允许轨道半径的量子化。
\paragraph{思考题 28.7}
按经典轨道模型，原子角动量没有量子化，
银原子通过Stern--Gerlach磁场的轨迹将出现什么分布？
为什么实验中出现离散分束的事实不能用经典轨道连续取向解释？
\paragraph{思考题 28.8}
基态氦原子两个电子的四个量子数
\(n,l,m_l,m_s\) 各是多少？为什么它们并不违反泡利原理？
\paragraph{思考题 28.9}
保持X光管电压不变，将银靶更换为铜靶，
X射线谱的短波截止和特征 \(K_\alpha\) 谱线波长分别如何变化？
\paragraph{思考题 28.10}
光子是费米子还是玻色子？它们是否受到Pauli不相容原理限制？
\paragraph{思考题 28.11}
什么是粒子数反转？为什么这种状态的两能级子系统可以用形式上的负热力学温度表征？
\paragraph{思考题 28.12}
在激光器出光端安装倾角等于布儒斯特角的玻璃窗（布儒斯特窗），
为什么透射激光可具有线偏振性质？其电矢量相对于入射面如何取向？
\paragraph{思考题 28.13}
电子能级与分子的振动、转动能级在能量标度上有什么不同？
分子光谱中的带状结构如何产生？
\paragraph{思考题 28.14}
为什么低温下分子转动和振动自由度可能分别“冻结”，
对定容摩尔热容不再产生经典均分定理所预期的贡献？
电子内部能级是否也可以热冻结？说明能隙与 \(k_BT\) 的关系。
